% Chapter 6, section 2: Admission, Batch, Plan, Apply
% Write in plain LaTeX. Code of a guide step: \guide{<step>}; commands: \begin{shell}...\end{shell};
% a screenshot: \screenshot{file}{caption}; a gap in the draft: \placeholder{...}.
% See book/latex/queen.sty for the full list.

\section{Admission, Batch, Plan, Apply}

A write goes through four steps.

\queenfig{write-path}{The four steps of a write.}

\paragraph{Admission.} Every node has a memory budget for the writes waiting to be planned (128 MiB
by default). A push takes its share before its body is read. When the budget is full, the push
waits, so a fast producer just slows down. After about 15 seconds it gets a 429, and the client
tries again. Pops and acks never wait here. We added this after a broker filled 10 GB of memory
with waiting requests and was killed.

\paragraph{Batch.} The leader collects up to 4,096 commands (or 4 MiB) and plans them together, as
one log entry. Up to eight entries are in flight at the same time. A retried request is found by
its id and gets the answer it missed, so nothing is planned twice.

\paragraph{Plan.} The planner decides what each command means right now, and writes the answer
down as facts: ``append these messages to partition 17, at offsets 4,210 to 4,309, at this
microsecond''. Ids, offsets and times are chosen here, once. To use more cores, eight lanes plan in
parallel: each lane owns a share of the partitions, and one last step plans what touches more than
one lane.

\queenfig{lanes}{Eight lanes plan in parallel; one step plans what crosses lanes.}

\paragraph{Apply.} Every node executes the committed entries, in order. Apply reads no clock and no
random number, so every node ends with the same bytes. Four threads apply the work of different
partitions at the same time.
